\documentclass[answers,12pt,addpoints]{exam} 
\usepackage{import}

\import{C:/Users/prana/OneDrive/Desktop/MathNotes}{style.tex}

% Header
\newcommand{\name}{Pranav Tikkawar}
\newcommand{\course}{Spatial Stats}
\newcommand{\assignment}{Homework 6.3}
\author{\name}
\title{\course \ - \assignment}

\begin{document}
\maketitle

\begin{questions}
\question 7.5


\begin{solution}
For the power law model with generalized autocovariance
$K_\phi(r) = \phi_1 \Gamma(-\phi_2) r^{2\phi_2}$ on the $5\times 5$
grid, with 24 linearly independent contrasts as differences
between the center observation at $(0.5,0.5)$ and each of the other
24 locations.  Let $D$ denote the $25\times 24$ matrix whose columns
are these contrast vectors, and let $C(\phi_2)$ be the $25\times 25$
generalized covariance matrix with entries
$C_{ij}(\phi_2) = K_\phi(\lVert s_i-s_j\rVert)$.  The covariance matrix
of the contrasts is then
\[
\Sigma(\phi_2) = D^\top C(\phi_2) D,
\]
which is positive definite of dimension $24\times 24$.

I computed the eigenvalues of $\Sigma(\phi_2)$ numerically for
$\phi_2 = 1 - 0.1^k$ with $k=2,3,4,5$, i.e.\ for
$\phi_2 \in \{0.99, 0.999, 0.9999, 0.99999\}$.  The largest five
eigenvalues are reported in Table~\ref{tab:eigs_top5}.
\end{solution}

\begin{table}[ht]
\centering
\begin{tabular}{cccccc}
\hline
$\phi_2$ & $\lambda_{1}$ & $\lambda_{2}$ & $\lambda_{3}$ & $\lambda_{4}$ & $\lambda_{5}$ \\
\hline
0.99    & 625.296408   & 625.296408   & 12.171594 & 2.560734 & 1.355417 \\
0.999   & 6250.232117  & 6250.232117  & 12.069184 & 2.556182 & 1.348082 \\
0.9999  & 62500.225674 & 62500.225674 & 12.059151 & 2.555755 & 1.347365 \\
0.99999 & 625000.200000& 625000.200000& 12.058150 & 2.555712 & 1.347294 \\
\hline
\end{tabular}
\caption{Largest eigenvalues $\lambda_1,\dots,\lambda_5$ of the
$24\times24$ contrast covariance matrix $\Sigma(\phi_2)$ for selected
values of $\phi_2$.}
\label{tab:eigs_top5}
\end{table}

\begin{solution}
In all cases the largest two eigenvalues are equal, and they increase
by approximately a factor of ten each time $\phi_2$ moves closer to 1.
Specifically, for $\phi_2 = 0.99$ I obtained
$\lambda_1 = \lambda_2 \approx 6.25\times 10^2$, while for
$\phi_2 = 0.99999$ these grow to
$\lambda_1 = \lambda_2 \approx 6.25\times 10^5$, with the remaining
22 eigenvalues staying of order one.  Thus, as $\phi_2 \to 1$ the
eigenstructure becomes increasingly dominated by a two-dimensional
leading eigenspace: almost all of the variance of the 24 contrasts is
concentrated in two directions, while variation in the other 22
directions is comparatively negligible.

To interpret these dominant directions, I mapped the associated
eigenvectors back onto the $5\times5$ grid (inserting a zero at the
center location) and plotted them as images.  For each $\phi_2$ the
first eigenvector shows a smooth gradient across the grid, increasing
from one corner to the opposite corner, whereas the second eigenvector
shows a similar gradient along the orthogonal diagonal.  These
patterns are essentially unchanged as $\phi_2$ varies from $0.99$ to
$0.99999$, indicating that the \emph{shape} of the leading spatial
modes is stable; only their eigenvalues change.  Together with
Table~\ref{tab:eigs_top5}, this shows that near the boundary
$\phi_2 \approx 1$ the power law model allocates almost all of the
variance of the intrinsic contrasts to two large-scale gradient
patterns over the grid, while higher-frequency or more localized
contrast patterns (associated with the smaller eigenvalues) contribute
relatively little.  This behaviour reflects the strong long-range
dependence of the power law covariance as $\phi_2$ approaches~1.

\begin{center}
    \includegraphics[scale = .3]{img/2eigvec.png}
\end{center}
\end{solution}

\begin{remark}
    \textbf{Excerise 4}

\begin{solution}
For a Gaussian process with power law generalized
autocovariance and true value of $\phi_2$ substantially less than $1$,
the profile restricted log-likelihood difference
$pr\ell(0.99999)-pr\ell(0.9999)$ will tend to change very slowly as the
number of observations $n$ increases.  A heuristic explanation can be
given by examining how the covariance matrix depends on $\phi_2$ near
the boundary point $\phi_2 = 1$.

The generalized autocovariance has the form
$K_\phi(r) = \phi_1 \Gamma(-\phi_2) r^{2\phi_2}$, so for any fixed
distance $r>0$ we have
\[
K_\phi(r; \phi_2+\delta)
  = K_\phi(r; \phi_2)\,\exp\!\bigl(2\delta \log r + o(\delta)\bigr)
\]
as $\delta \to 0$.  When $\phi_2$ is already extremely close to $1$,
moving from $\phi_2 = 0.9999$ to $\phi_2 = 0.99999$ corresponds to a
tiny increment $\delta = 9\times 10^{-5}$, so each individual
covariance entry changes only by a relative amount of order
$\delta \log r$, which is very small for the fixed interpoint distances
in a typical spatial sampling design.  In matrix terms, the covariance
matrices $C(0.9999)$ and $C(0.99999)$ differ by a small perturbation in
operator norm.

For Gaussian data, the (restricted) log-likelihood can be written as
\[
pr\ell(\phi_2)
= -\tfrac{1}{2}\Bigl\{\log\det C(\phi_2)
   + Z^\top C(\phi_2)^{-1} Z\Bigr\} + \text{const},
\]
after profiling out the mean and scale parameters.  Since both
$\log\det C(\phi_2)$ and $C(\phi_2)^{-1}$ are smooth functions of
$C(\phi_2)$, a small perturbation of size $\delta$ in $C$ produces a
change in $pr\ell(\phi_2)$ of order $\delta^2$ in a neighbourhood where
the likelihood surface is reasonably flat.  Thus the difference
$pr\ell(0.99999) - pr\ell(0.9999)$ is itself of very small order in
$\delta$.

When the true value of $\phi_2$ is substantially less than $1$, both
$0.9999$ and $0.99999$ lie far out in the tail of the likelihood, in a
region where the data provide little information to distinguish between
them.  Increasing the sample size $n$ refines the estimate of $\phi_2$
in the interior of the parameter space but does not appreciably change
the relative support for these two extreme boundary values, since they
induce almost indistinguishable covariance structures.  Consequently
the vertical gap $pr\ell(0.99999)-pr\ell(0.9999)$ changes only very
slowly as $n$ grows. 
\end{solution}


\end{remark}




\question 7.6

\begin{solution}
I compared the expected Fisher
information for REML estimation under a two-parameter power law
generalized covariance and a two-parameter exponential covariance,
each with an unknown constant mean.  For each model I fixed a
parameter value and computed the $2\times2$ expected Fisher
information matrix $I(\theta)$ for four different designs of $n=25$
observation locations in the unit square:
\begin{itemize}
\item a regular $5\times5$ grid (\texttt{grid});
\item a slightly jittered version of the grid (\texttt{jitter});
\item 25 points sampled independently from $\operatorname{Unif}([0,1]^2)$
(\texttt{uniform});
\item a central cluster with points sampled from a bivariate normal
centered at $(0.5,0.5)$ and truncated to $[0,1]^2$ (\texttt{center}).
\end{itemize}
For the power law model I used $\phi_1=1$, $\phi_2=0.7$, and for the
exponential model I used $\sigma^2=1$, $\rho=0.3$.

For each design and model I formed the covariance matrix
$\Sigma(\theta)$ and its partial derivatives with respect to the two
covariance parameters.  With design matrix $X=\mathbf{1}_n$, the
expected REML Fisher information for parameters $\theta_a,\theta_b$ is
\[
I_{ab}
= \tfrac{1}{2} \operatorname{tr}\!\left\{
P\,\Sigma_{,a}\,P\,\Sigma_{,b}\right\},\qquad
P = \Sigma^{-1} - \Sigma^{-1} X
       (X^\top \Sigma^{-1} X)^{-1} X^\top \Sigma^{-1},
\]
which I evaluated numerically for each design.
To summarize the information matrices I recorded, for each
$2\times2$ matrix $I(\theta)$, the trace
$\operatorname{tr}(I)$, determinant $\det(I)$, and condition number
$\kappa(I)$ (ratio of largest to smallest eigenvalue).  The results
are shown in Table~\ref{tab:fisher_power} for the power law model and
Table~\ref{tab:fisher_exp} for the exponential model.
\end{solution}

\begin{table}[ht]
\centering
\begin{tabular}{lccc}
\hline
Design & $\operatorname{tr}(I)$ & $\det(I)$ &
$\kappa(I)$ \\
\hline
grid    & 484.0 & $1.46\times 10^{4}$ & 13.99 \\
jitter  & 502.8 & $1.54\times 10^{4}$ & 14.36 \\
uniform & 627.2 & $2.51\times 10^{4}$ & 13.63 \\
center  & 524.5 & $1.69\times 10^{4}$ & 14.23 \\
\hline
\end{tabular}
\caption{Trace, determinant, and condition number of the Fisher
information matrix for the power law covariance model under four
sampling designs.}
\label{tab:fisher_power}
\end{table}

\begin{table}[ht]
\centering
\begin{tabular}{lccc}
\hline
Design & $\operatorname{tr}(I)$ & $\det(I)$ &
$\kappa(I)$ \\
\hline
grid    & 80.1  & $3.25\times 10^{2}$ & 17.69 \\
jitter  & 80.7  & $3.34\times 10^{2}$ & 17.48 \\
uniform & 98.3  & $3.28\times 10^{2}$ & 27.43 \\
center  & 130.8 & $1.79\times 10^{2}$ & 93.31 \\
\hline
\end{tabular}
\caption{Trace, determinant, and condition number of the Fisher
information matrix for the exponential covariance model under the same
four sampling designs.}
\label{tab:fisher_exp}
\end{table}

\begin{solution}
    

For the power law model, both the trace and determinant of $I(\theta)$
vary substantially across designs.  The trace increases from about
$4.84\times10^{2}$ for the regular grid to $6.27\times10^{2}$ for the
uniform design, while the determinant increases from about
$1.46\times10^{4}$ to $2.51\times10^{4}$, indicating a marked gain in
both marginal and joint precision of $(\phi_1,\phi_2)$ when the design
is changed.The central cluster design also differs
appreciably from the grid, with intermediate trace and determinant
values.  In contrast, the condition numbers remain in a relatively
narrow range (about 13.6--14.4), so the correlation structure between
the two parameters does not change dramatically with design; the main
effect is a sizeable change in the overall amount of information.

For the exponential model the picture is noticeably different.  The
trace increases from about $8.0\times10^{1}$ for the grid and jittered
designs to $9.8\times10^{1}$ for the uniform design and
$1.31\times10^{2}$ for the center design, a smaller relative change
than in the power law case, the determinant
stays close to $3.3\times10^{2}$ for the grid, jittered, and uniform
designs, and only drops substantially for the highly clustered center
design, indicating that the joint precision of $(\sigma^2,\rho)$ is
comparatively stable across most designs. While the
condition number does become large for the center design, reflecting
stronger correlation between the exponential parameters, the overall
variation in $\operatorname{tr}(I)$ and $\det(I)$ across designs is
smaller than for the power law model.

Taken together, these numerical experiments support the claim: the statistical properties of REML estimates (as reflected in
the Fisher information) are much more sensitive to the choice of
observation locations for the two-parameter power law covariance model
than for the corresponding two-parameter exponential model.
\end{solution}

\end{questions}

\end{document}